\chapter{Projectors and quantum packets}
\label{ch:projectors}
\index{projector}
\index{packet decomposition}

\section{Orthogonal packet decompositions}
Let $\Hh$ be a finite-dimensional Hilbert space and let $\{P_\alpha\}_{\alpha\in A}$ be orthogonal projectors satisfying
\[
P_\alpha P_\beta=\delta_{\alpha\beta}P_\alpha,
\qquad
\sum_{\alpha\in A}P_\alpha=I.
\]
Every state and operator decompose as
\[
\ket\psi=\sum_\alpha P_\alpha\ket\psi,
\qquad
X=\sum_{\alpha,\beta}P_\alpha X P_\beta.
\]
The diagonal blocks act within packets. The off-diagonal blocks transfer amplitude or coherence between packets.

Define the packet-transfer graph of an operator $X$ by placing an edge $\beta\to\alpha$ whenever
\[
P_\alpha X P_\beta\neq0.
\]
An exactly symmetry-preserving operator has only diagonal edges. A sparse transfer graph may still permit efficient simulation if the active packets and internal states remain compact.

\section{Conserved observables}
Let
\[
Q=\sum_k q_kP_k
\]
be a self-adjoint observable with spectral projectors $P_k$.

\begin{theorem}[Invariant quantum packets]
\label{thm:invariant-packets}
If a unitary $U$ satisfies $[U,Q]=0$, then
\[
P_iUP_j=0\qquad(i\neq j),
\]
and
\[
U=\bigoplus_i U_i,
\qquad
U_i=P_iUP_i\big|_{P_i\Hh}.
\]
If $P_k\ket\psi=\ket\psi$, then $P_kU\ket\psi=U\ket\psi$.
\end{theorem}

\begin{proof}
Since $U$ commutes with $Q$, it commutes with every spectral projector obtained from the functional calculus of $Q$. Thus
\[
P_iUP_j=UP_iP_j=0
\]
for $i\neq j$. Summing the diagonal blocks gives the direct-sum form. The final statement follows from
\[
P_kU\ket\psi=UP_k\ket\psi=U\ket\psi.
\]
\end{proof}

\section{Residue and character packets}
Suppose $Q$ has integer spectrum. For a modulus $M$, define
\[
B_r^{(M)}=
\sum_{k\equiv r\pmod M}P_k.
\]
Let
\[
V_M=\exp\left(\frac{2\pi i}{M}Q\right),
\qquad
\omega_M=e^{2\pi i/M}.
\]
Then
\begin{equation}
B_r^{(M)}
=
\frac1M\sum_{\ell=0}^{M-1}
\omega_M^{-r\ell}V_M^\ell.
\label{eq:residue-projector}
\end{equation}
On a $Q$-eigenvector, the right side is the roots-of-unity filter. The same finite-character formula selects residue classes in signal polyphase analysis and charge sectors in quantum theory.

\section{Finite-group projectors}
Let a finite group $G$ act unitarily by $U_g$. For an irreducible character $\chi$ of degree $d_\chi$, the central isotypic projector is
\begin{equation}
P_\chi
=
\frac{d_\chi}{|G|}
\sum_{g\in G}
\overline{\chi(g)}U_g.
\label{eq:character-projector}
\end{equation}
Character orthogonality gives
\[
P_\chi P_\psi=\delta_{\chi\psi}P_\chi,
\qquad
\sum_\chi P_\chi=I
\]
on the represented space. If an operator $A$ commutes with every $U_g$, then
\[
P_\chi A P_\psi=0
\qquad(\chi\neq\psi).
\]

This framework includes momentum, parity, total spin, point-group symmetry, permutation symmetry, and stabilizer syndrome labels. The isotypic projector identifies a symmetry type. It does not necessarily resolve the multiplicity space within that type.

\section{Nested refinements}
If one conserved label is too coarse, intersect commuting projector families. For mutually commuting observables $Q_1,\ldots,Q_s$, a joint packet has the form
\[
P_{\bm\alpha}
=
P^{(1)}_{\alpha_1}\cdots P^{(s)}_{\alpha_s}.
\]
The dimension can decrease dramatically, but each extra label must be efficiently computable and preserved. A refinement that requires exponential preprocessing merely relocates the complexity.

\begin{remark}
A packet decomposition is basis independent as an orthogonal direct sum, although a practical simulator must choose a concrete basis or compressed coordinate system within each active packet.
\end{remark}
